\documentclass[12pt,a4paper]{ctexart}
\usepackage[left=2.5cm,right=2.5cm,top=2.5cm,bottom=2.5cm]{geometry}
\usepackage{amsmath}
\usepackage{array}
\usepackage{booktabs}
\usepackage{graphicx}
\usepackage{hyperref}
\usepackage{longtable}
\usepackage{setspace}
\usepackage{caption}
\captionsetup{font=small}
\setstretch{1.25}

\title{基于U-Net及其变体的遥感影像道路提取实验报告}
\author{}
\date{2026年3月15日}

\begin{document}
\maketitle

\begin{abstract}
本报告围绕遥感影像道路提取任务，基于 Massachusetts Roads Dataset 重建了完整的 PyTorch 训练、评估与整图推理流程，并以 U-Net 为基线开展多组结构改进实验。研究最初计划以集成轻量级双线性注意力模块的 DLGU-Net 作为主要创新模型，但实际实验结果表明，DLGU-Net 在统一整图评估协议下未能超过基线模型。统一对比结果显示，基线 U-Net 的测试集 IoU 为 0.645964，F1 为 0.782378；在保持基线主体结构不变的前提下，于瓶颈层引入空洞卷积上下文模块后，测试集 IoU 提升至 0.646839，F1 提升至 0.782933，取得了小幅增益。除此之外，实验还额外实现并测试了 Attention U-Net、UNet++、DDU-Net、预训练 ResNet34 编码器 U-Net、原始风格 Vanilla U-Net 及其残差版本。补充实验中，Vanilla U-Net 取得最高测试集 IoU 0.647187，但由于其同时调整了归一化方式、上采样方式与通道规模，不能视为对基线的严格等价改进。综合来看，本工作的主要结论是：在当前数据集与训练条件下，基线 U-Net 仍然是最稳健的参考模型；面向基线的直接结构增强中，空洞卷积上下文模块表现出有限但可量化的正向作用；更复杂的注意力、多解码器与预训练编码器结构未表现出稳定优势。
\end{abstract}

\tableofcontents
\newpage

\section{选题的背景与意义}
随着高分辨率对地观测技术的快速发展，遥感影像已经成为城市规划、地图更新、自动驾驶高精地图构建、灾害应急响应与军事侦察的重要基础数据。道路作为地表最核心的人工线性目标之一，其提取质量直接影响后续的网络分析、路径规划、地图矢量化和交通监测精度。相比人工判读，基于深度学习的自动道路提取具有处理速度快、可扩展性强和适合大规模生产等优势，因此具有明确的研究价值与应用价值。

遥感影像道路提取的困难主要体现在三个方面。第一，道路目标经常受到建筑物阴影、植被遮挡、停车场与屋顶纹理干扰，导致前景与背景混淆。第二，道路在几何上具有细长、连续、分支丰富和拓扑结构复杂等特点，模型不仅要识别局部纹理，还要维护连通性与方向一致性。第三，公开道路数据集规模有限，模型性能容易受数据分布、采样规模和训练策略影响，导致复杂模型未必能稳定优于简单模型。

在此背景下，以 U-Net 为代表的编码器--解码器结构仍然是道路提取任务的重要基线。该类模型结构清晰、训练稳定、易于扩展，适合作为各类结构增强方法的对比参照。本研究的意义不在于堆叠复杂模块，而在于以统一的数据处理、训练配置和整图评估协议为基础，验证不同结构改动是否能够带来真实、可复现的性能提升。最终结果表明，数据导向的验证比结构设想本身更重要，这也是本报告与开题阶段相比最关键的调整。

\section{设计的思路与主要内容}
本研究以 U-Net 为主线展开。开题阶段的原始思路是围绕道路断裂与背景干扰问题，在 U-Net 的跳跃连接处嵌入轻量级双线性注意力模块 DLAM，构建 DLGU-Net，从而增强模型对道路关键特征的聚焦能力并提升细长道路的连通性。

在实际实现过程中，工作内容扩展为三个层次。第一层是工程重建。由于原始项目代码存在缺失，重新构建了数据整理、离线 patch 生成、模型注册、训练器、整图评估与单图推理流程，使全部实验能够在同一代码框架下复现。第二层是主线模型验证。围绕基线 U-Net 与 DLGU-Net，统一完成训练、验证和整图测试，并据此判断开题方案是否成立。第三层是补充探索。为了判断性能差异究竟来自注意力机制还是来自主干结构本身，进一步实现并测试了 Attention U-Net、UNet++、DDU-Net、预训练 ResNet34 编码器 U-Net、Vanilla U-Net、残差 Vanilla U-Net 以及带空洞卷积上下文模块的基线变体。

因此，本报告的主要内容不再沿用“DLGU-Net 优于所有方法”的预设结论，而是基于实验数据重新组织为以下三个问题：基线模型能达到什么水平；DLAM 注意力模块是否真的有效；哪些结构变化在当前数据集上能够取得稳定的增益。

\section{毕业设计所用的方法与技术路线}
\subsection{数据集与预处理}
本研究的全部定量实验均基于 Massachusetts Roads Dataset 完成。当前本地实验数据共包含训练集 1108 张、验证集 14 张、测试集 49 张遥感影像及其对应标注。原始图像尺寸约为 $1500 \times 1500$。为保证网络输入尺寸统一，先将图像与标签补边至 $1536 \times 1536$，再按 $3 \times 3$ 网格切分为 9 个 $512 \times 512$ patch。训练集和验证集采用离线 patch 数据集以提高吞吐效率，测试集则采用整图切块推理与拼接恢复的评估流程。

图像预处理采用 ImageNet 均值和方差归一化，标签采用二值化处理。训练过程中不额外引入复杂数据增强，以减少模型之间由于增强策略差异带来的比较偏差。与最初在线整图读取并实时切 patch 的实现相比，离线 patch 方案显著降低了 CPU 数据准备开销，并将基线单 epoch 的训练时间从早期实现的约 409 秒压缩到约 119 秒量级，从而保证了后续大规模对比实验可以完成。

\subsection{基线模型}
基线模型采用四层编码器--解码器结构。卷积块由两层 $3 \times 3$ 卷积、BatchNorm 和 ReLU 构成；下采样使用 MaxPool；上采样采用双线性插值后与编码端特征拼接，再通过卷积块完成特征融合。模型初始通道数设置为 32，总参数量为 7.85M。该结构对应当前项目中的标准比较对象，也是全文中“baseline”的唯一指代。

\subsection{DLGU-Net 与 DLAM 模块}
DLGU-Net 在基线 U-Net 的四个跳跃连接分支上插入 DLAM 模块。DLAM 由通道注意力与空间注意力两部分组成。通道注意力分支对输入特征分别进行全局平均池化和全局最大池化，并通过共享 MLP 生成通道权重；空间注意力分支使用条形卷积核提取水平方向与垂直方向的长程上下文信息，生成空间权重图。两者相乘后与输入特征做残差融合，表示为
\[
F_{\text{out}} = F + F \odot A_c(F) \odot A_s(F),
\]
其中 $A_c(\cdot)$ 表示通道注意力，$A_s(\cdot)$ 表示空间注意力，$\odot$ 表示逐元素乘法。

该设计的出发点是合理的：道路具有明显的线性延展结构，通道重标定有助于强化道路语义，空间重标定有助于突出线状响应。然而，结构上的理论动机并不等价于最终性能增益，因此需要通过完整实验进行验证。

\subsection{补充模型与对照模型}
为了从更宽的范围理解结构改动的有效性，本研究还实现了多种补充模型：
\begin{enumerate}
    \item Attention U-Net：在跳跃连接前引入门控注意力。
    \item UNet++：在保持基线卷积块和上采样方式的前提下，引入嵌套跳跃连接。
    \item DDU-Net：采用 ResNet-50 编码器、双解码器和上下文模块。
    \item ResNet34 U-Net：以预训练 ResNet34 作为编码器。
    \item Vanilla U-Net：更接近原始 U-Net 结构，不使用 BatchNorm，采用转置卷积上采样，初始通道数为 64。
    \item Residual Vanilla U-Net：在 Vanilla U-Net 的双卷积块中加入残差连接。
    \item Dilated Baseline U-Net：在基线 U-Net 的瓶颈层后加入并行空洞卷积上下文模块。
\end{enumerate}

其中，Dilated Baseline U-Net 是本文最重要的“基于 baseline 的直接改进”模型。其核心结构写为
\[
F_{\text{ctx}} = F + \phi([D_1(F), D_2(F), D_4(F), D_8(F)]),
\]
其中 $D_r(\cdot)$ 表示膨胀率为 $r$ 的空洞卷积，$[\cdot]$ 表示通道拼接，$\phi(\cdot)$ 表示 $1 \times 1$ 卷积融合。该结构的目的是在不改变整体编码器--解码器骨架的前提下，增强瓶颈层的多尺度上下文表达。

\subsection{训练策略与评估协议}
训练框架基于 PyTorch 实现，硬件环境为 NVIDIA GeForce RTX 5090。基线、DLGU-Net、Attention U-Net、UNet++ 和 Dilated Baseline U-Net 等主对比模型统一采用 batch size 为 4、初始学习率为 $1 \times 10^{-4}$、Dice+BCE 组合损失、StepLR 学习率调度和早停策略。训练阶段使用 patch 级数据加载；最终模型比较统一采用整图评估协议，即对完整图像进行切块预测、拼接恢复后，再计算 IoU、F1、Precision 与 Recall。

需要强调的是，训练日志中的 patch 验证指标与最终整图验证指标并非同一统计对象。因此，本报告的模型排序与结论全部以统一整图验证和整图测试结果为准，而不以训练过程中记录的 patch 级验证曲线为准。

\section{实验设置与结果统计口径}
为了保证比较公平，本文采用以下统一原则：
\begin{enumerate}
    \item 数据集统一为 Massachusetts Roads Dataset。
    \item 最终指标统一采用整图验证集与整图测试集结果。
    \item 排序主要依据测试集 IoU，其次参考测试集 F1。
    \item 训练失败、中断实验和 smoke 实验不纳入结果表。
\end{enumerate}

开题阶段计划中的 DeepGlobe 跨数据集泛化实验，在当前报告中未纳入最终量化对比表。一方面，当前本地完整、统一、可复核的定量结果集中在 Massachusetts 数据集；另一方面，为避免不同数据准备状态带来的额外变量，本文最终只报告完全在同一流程下得到的实验结果。

\section{实验结果与分析}
\subsection{全部模型的统一结果}
表 \ref{tab:all-results} 给出了当前全部完整实验模型在统一整图评估协议下的结果。为避免歧义，表中验证集结果与测试集结果均来源于整图评估，而非训练阶段的 patch 级验证。

\begin{table}[htbp]
\centering
\caption{全部模型整图验证与整图测试结果对比}
\label{tab:all-results}
\resizebox{\textwidth}{!}{
\begin{tabular}{lcccccccc}
\toprule
模型 & 参数量(M) & Val IoU & Val F1 & Test IoU & Test F1 & Test Precision & Test Recall & 说明 \\
\midrule
Vanilla U-Net & 31.03 & \textbf{0.6321} & \textbf{0.7722} & \textbf{0.6472} & \textbf{0.7832} & 0.7954 & 0.7748 & 补充探索，非严格等价基线改进 \\
Dilated Baseline U-Net & 18.34 & 0.6273 & 0.7684 & 0.6468 & 0.7829 & \textbf{0.8111} & 0.7593 & 基于 baseline 的直接结构增强 \\
Baseline U-Net & 7.85 & 0.6298 & 0.7703 & 0.6460 & 0.7824 & 0.7877 & \textbf{0.7802} & 全文基线模型 \\
UNet++ & 8.30 & 0.6308 & 0.7711 & 0.6432 & 0.7802 & 0.7910 & 0.7725 & 基于 baseline 的嵌套跳跃连接 \\
DLGU-Net & 8.30 & 0.6261 & 0.7675 & 0.6419 & 0.7792 & 0.7932 & 0.7694 & 基于 DLAM 的主线模型 \\
Residual Vanilla U-Net & 32.43 & 0.6306 & 0.7710 & 0.6389 & 0.7768 & 0.8116 & 0.7486 & Vanilla U-Net 残差化 \\
Attention U-Net & 7.98 & 0.6268 & 0.7680 & 0.6369 & 0.7755 & 0.7982 & 0.7576 & 跳跃连接门控注意力 \\
ResNet34 U-Net & 29.17 & 0.6222 & 0.7646 & 0.6284 & 0.7691 & 0.7748 & 0.7666 & 预训练编码器 \\
DDU-Net & 46.78 & 0.6163 & 0.7602 & 0.6277 & 0.7693 & 0.7322 & 0.8138 & 双解码器与更重上下文模块 \\
\bottomrule
\end{tabular}
}
\end{table}

\subsection{基线模型与 DLGU-Net 的对比}
从最终结果看，DLGU-Net 未能实现开题阶段预期。基线 U-Net 的测试集 IoU 为 0.645964，DLGU-Net 为 0.641934，下降 0.004030；基线 F1 为 0.782378，DLGU-Net 为 0.779244，下降 0.003134。进一步观察 Precision 与 Recall 可以发现，DLGU-Net 的 Precision 从 0.787734 上升到 0.793195，但 Recall 从 0.780151 下降到 0.769380，说明注意力模块使模型预测更保守，抑制了部分背景误报，但同时削弱了道路召回，最终没有转化为更好的 IoU 和 F1。

这一定量结果意味着：DLAM 作为结构设想是可实现的，但在当前数据集、训练配置和 U-Net 主干条件下，并未形成稳定增益。因此，报告在结论表达上必须与开题阶段区分开来。开题阶段强调的是“拟提出的改进方向”，而当前实验支持的结论是“DLGU-Net 已完成设计与验证，但其性能未超过基线模型”。

\subsection{基于 baseline 的直接改进：空洞卷积上下文模块}
在所有基于 baseline 的直接结构增强中，Dilated Baseline U-Net 的结果最接近“对 baseline 的小幅改进”。该模型仅在瓶颈层后加入并行空洞卷积上下文模块，其余主体结构保持不变。实验结果表明，其测试集 IoU 为 0.646839，相较 baseline 提升 0.000875；测试集 F1 为 0.782933，相较 baseline 提升 0.000555。

不过，这一增益并不完全稳定。整图验证集上，Dilated Baseline U-Net 的 IoU 为 0.627253，低于 baseline 的 0.629811。换言之，空洞卷积模块在测试集上取得了极小的正向提升，但验证集表现未同步提升。因此，本文对该结果的表述应限定为“取得了小幅测试增益”，而不能表述为“稳定且显著优于 baseline”。

\subsection{Vanilla U-Net 的补充说明}
Vanilla U-Net 取得了当前全部实验中的最高测试集 IoU 0.647187 和最高测试集 F1 0.783246，相比 baseline 分别提升 0.001223 和 0.000868。该结果说明，在当前任务中，去除 BatchNorm、采用转置卷积上采样并扩大通道规模后，主干网络仍然具有进一步提升空间。

但是，Vanilla U-Net 不能直接归入“对 baseline 的严格等价改进”。原因在于该模型同时改变了三项关键因素：卷积块归一化方式、解码器上采样方式以及初始通道数。其参数量从 baseline 的 7.85M 增加到 31.03M，比较关系已经不再是单一变量控制实验。因此，Vanilla U-Net 更适合作为补充探索结果，用于说明本研究在主线模型之外完成了大量结构对照工作，而不宜被表述为本文对 baseline 的唯一主结论。

\subsection{复杂模型不一定更优}
UNet++、Attention U-Net、DDU-Net、ResNet34 U-Net 和 Residual Vanilla U-Net 的结果共同说明：在当前数据规模和任务条件下，更复杂的结构并不会自然带来更好的道路提取性能。尤其是 DDU-Net，虽然参数量达到 46.78M，且 Recall 达到最高的 0.8138，但 Precision 仅为 0.7322，最终 IoU 反而降至 0.6277。类似地，ResNet34 U-Net 和 Attention U-Net 也均未超过基线。

这一现象说明，当前任务更依赖结构稳定性与前景/背景平衡控制，而不是单纯扩大模型容量或叠加控制分支。换言之，遥感道路提取任务对结构复杂度的收益并不线性，复杂设计可能只是改变 Precision 与 Recall 的分配方式，而不一定减少总体误差。

\section{工作量与完成情况}
从工程实现角度看，本研究的工作量主要体现在以下方面：
\begin{enumerate}
    \item 重建完整的道路提取实验工程，包括数据准备、离线 patch 生成、模型构建、训练器、整图评估与推理脚本。
    \item 将 Massachusetts 数据集规范化整理为统一目录结构，并构建离线 patch 数据集，提高训练吞吐效率。
    \item 完整实现并训练了 9 个可对比模型，覆盖基线模型、注意力模型、嵌套跳跃连接模型、双解码器模型、预训练编码器模型、原始风格 U-Net 及其变体。
    \item 对全部完整模型重新执行统一整图验证与整图测试，消除了 patch 验证与整图测试在统计口径上的差异。
    \item 对开题主线模型 DLGU-Net 进行了真实验证，并对不符合预期的结果重新组织叙述，使报告结论与实验事实一致。
\end{enumerate}

因此，本工作的价值不仅体现在单一模型指标上，还体现在构建了一个可复核、可扩展、能够支持多模型公平比较的实验框架。该部分工作为后续继续开展消融实验和结构微调提供了稳定基础。

\section{结论}
基于当前统一整图评估结果，可以得出以下结论。

第一，Baseline U-Net 在 Massachusetts Roads Dataset 上表现稳健，是当前实验体系中的核心参照模型。第二，开题阶段提出的 DLGU-Net 已经完成设计、实现和系统验证，但其在当前数据集上的性能未能超过 baseline，因此不能将其表述为最终最优模型。第三，在基于 baseline 的直接结构增强中，空洞卷积上下文模块带来了极小的测试集性能提升，说明该方向具有一定有效性，但提升幅度有限且验证集一致性不足。第四，Vanilla U-Net 取得了全部模型中的最优测试结果，但由于其改变了主干结构与参数规模，更适合作为补充探索成果，而不是对 baseline 的严格单变量改进。

综合而言，本研究最终形成的结论是：遥感道路提取任务中的模型改进应当坚持数据导向与统一评估原则。以 baseline 为锚点开展的小幅结构增强，比直接堆叠复杂模块更容易获得可信结果；而所有模型结论都必须以统一整图评估结果为依据，而不能仅依据训练过程中的局部验证指标。

\end{document}
